\documentclass[10pt,a4paper]{amsart}
\usepackage[margin=19mm]{geometry}
\usepackage{amsmath,amssymb,mathtools}
\usepackage[scr=rsfs,cal=euler]{mathalfa}
\usepackage{xfrac}
\usepackage[unicode,hidelinks,bookmarks=false]{hyperref}
\newcommand{\ie}{i.e.\ }
\newcommand{\cf}{cf.\ }
\newcommand{\resp}{respectively\ }
\newcommand{\Subsection}{\subsection}
\providecommand{\nobreakdash}{\nobreak\hbox{-}}
\setlength{\emergencystretch}{2em}
\multlinegap0pt
\usepackage{bm}
\usepackage{amssymb}
\newtheorem{lemma}{Lemma}
\newtheorem{theorem}{Theorem}
\newcommand{\F}{\mathbb{F}}
\DeclareMathOperator{\Span}{span}
\title{The Sidon Decomposition Problem in Abelian Groups of Bounded Torsion}
\author{Mark Lewko}
\date{Source: arXiv:2606.06669v1 (CC BY 4.0)}
\begin{document}
\maketitle

\noindent\emph{Source and licence.} The Sidon Decomposition Problem in Abelian Groups of Bounded Torsion. Version arXiv:2606.06669v1 (CC BY 4.0), distributed by arXiv under the Creative Commons Attribution 4.0 International licence, http://creativecommons.org/licenses/by/4.0/, as recorded in the arXiv licence metadata for this exact version and shown on the abstract page https://arxiv.org/abs/2606.06669v1. Full licence notice: \texttt{licenses/arxiv-2606.06669-CC-BY-4.0.txt}.\par\smallskip

\noindent\textit{Editorial scope.} Counted unit: the article's theorem thm:supportpartition (source0.tex lines 307--319). The article's complete original proof of it is delivered with it (source0.tex lines 321--427, one \texttt{proof} environment of 107 lines). No mathematics is written, repaired or re-derived: the statement and the proof are the article's own lines, in the article's own notation, and no retained line is substituted or re-flowed.  Retained uncounted as the source-local context the proof needs: lines 211--225.  Nothing was omitted from any retained span, so the package carries no omission record.  No retained line contains a citation, so the wrapper carries no bibliography and no bibliography entry.  No retained line contains a figure environment, an image-inclusion command or drawing code, so no image is delivered and the package declares no asset dependency.  Editorial formatting only, in addition: an AMS \texttt{amsart} preamble that restates the article's own declarations and macros used by the retained lines, namely the environments \texttt{lemma}, \texttt{theorem} on separate counters, together with the standard packages those lines need.  The retained environments print in this document as Lemma 1, Theorem 1 in order of appearance, whereas the article prints the same items with the numbering of its own full document.  The complete native source of the article as distributed in the arXiv e-print for this exact version is bundled unchanged as \texttt{sources/202201/source0.tex}.
\begin{lemma}\label{lem:complement-radohorn}
Let $X$ be finite, let $(v_x)_{x\in X}$ be a finite indexed list of vectors
spanning a finite-dimensional vector space $W$ over $\F$, and let $k\geq 2$.
Suppose that for every proper subspace $L<W$,
$$
        |\{x\in X:v_x\notin L\}|
        \geq
        \frac{k}{k-1}\bigl(\dim W-\dim L\bigr).
$$
Then $X$ admits a partition $X=X_1\sqcup\cdots\sqcup X_k$ such that, for each
$j$,
$$
        \Span\{v_x:x\notin X_j\}=W.
$$
\end{lemma}

\begin{theorem}\label{thm:supportpartition}
Let $X$ be finite, let $\mathcal C\subseteq \F^X\setminus\{0\}$, and, for
$Y\subseteq X$, define
$$
        \lambda_{\mathcal C}(Y)
        =
        \dim\Span\{c\in\mathcal C:\operatorname{supp} c\subseteq Y\}.
$$
Let $k\geq 2$. If $\lambda_{\mathcal C}(Y) \leq \frac{k-1}{k}|Y|$ for every
$Y\subseteq X$, then $X$ admits a partition $X=X_1\sqcup\cdots\sqcup X_k$
such that no vector $c\in\mathcal C$ has support contained in a single
class $X_j$.
\end{theorem}

\begin{proof}
We argue by induction on $|X|$. If $\mathcal C$ is empty the condition is
vacuously satisfied and any partition works.

The proof uses Lemma~\ref{lem:complement-radohorn}, which partitions an indexed list
of vectors in some ambient space $W$ into $k$ classes such that each class's
complement spans $W$. We want to apply this to produce a partition of $X$ such
that every $\mathbf{c} \in \mathcal C$ has a support point outside each
class. We assign to each $x \in X$ a vector $\mathbf{m}_x \in \F^r$
such that $x \in \operatorname{supp} \mathbf{c}$ if and only if
$\boldsymbol{\beta}_{\mathbf{c}} \cdot \mathbf{m}_x \neq 0$, where
$\boldsymbol{\beta}_{\mathbf{c}} \in \F^r$ is the coefficient vector of
$\mathbf{c}$. Given such vectors, if the complement of each class spans $\F^r$,
then for any nonzero $\boldsymbol{\beta}_{\mathbf{c}}$ there is a point outside
each class not orthogonal to $\boldsymbol{\beta}_{\mathbf{c}}$, which is
precisely a support point of $\mathbf{c}$ outside that class.

Let $\mathbf{b}_1, \ldots, \mathbf{b}_r$ be a basis of $R = \Span\mathcal C
\leq \F^X$, and form the $r \times |X|$ matrix whose rows are
$\mathbf{b}_1, \ldots, \mathbf{b}_r$, with $(\mathbf{b}_i)_x \in \F$ denoting
the $x$-th coordinate of $\mathbf{b}_i$. Its column $x$ is the vector
$\mathbf{m}_x = ((\mathbf{b}_1)_x, \ldots, (\mathbf{b}_r)_x)^\top \in \F^r$. Every
$\mathbf{c} \in \mathcal C$ lies in $R$, so $\mathbf{c} = \sum_{i=1}^r
\beta_i \mathbf{b}_i$ for unique coefficients, and the coefficient vector
$\boldsymbol{\beta}_{\mathbf{c}} = (\beta_1, \ldots, \beta_r)^\top \in \F^r$
satisfies
$$
        \mathbf{c}_x = \boldsymbol{\beta}_{\mathbf{c}} \cdot \mathbf{m}_x
        \quad \text{for every } x \in X.
$$
Thus $x \in \operatorname{supp} \mathbf{c}$ if and only if
$\boldsymbol{\beta}_{\mathbf{c}} \cdot \mathbf{m}_x \neq 0$, as required.
Any $x$ with $\mathbf{m}_x = 0$ satisfies $\mathbf{c}_x = 0$ for every
$\mathbf{c} \in \mathcal C$ and may be placed in any class; we therefore
restrict attention to $X^+ = \{x \in X : \mathbf{m}_x \neq 0\}$. The
hypothesis applied to $Y = X^+$ gives $|X^+| \geq \frac{k}{k-1}r$, and the
columns $(\mathbf{m}_x)_{x \in X^+}$ span $\F^r$, since otherwise some
nontrivial linear combination of the $\mathbf{b}_i$ would vanish on all of $X$.

We would like to apply Lemma~\ref{lem:complement-radohorn} directly to
$(\mathbf{m}_x)_{x \in X^+}$ in $\F^r$, but the lemma's hypothesis may fail
if some proper subspace $L \subsetneq \F^r$ contains too many
columns $\mathbf{m}_x$. We start by looking for the smallest subspace $W \leq \F^r$ such that the
columns landing in $W$ satisfy the lemma's hypothesis within $W$. We will apply
Lemma~\ref{lem:complement-radohorn} on $W$. Vectors that vanish on $W$ are passed to the
induction step. For each subspace $U \leq \F^r$, we write $X_U = \{x \in X^+ :
\mathbf{m}_x \in U\}$.
Choose $W \leq \F^r$ of minimal positive dimension satisfying $|X_W| \geq
\frac{k}{k-1}\dim W$. Since $\F^r$ satisfies this by the bound on
$|X^+|$, such a $W$ must exist.

We verify the hypothesis of Lemma~\ref{lem:complement-radohorn} for the indexed list
$(\mathbf{m}_x)_{x \in X_W}$ in the space $W$. Since every $x \in X_W
\subseteq X^+$ has $\mathbf{m}_x \neq 0$, no column lies in the zero subspace,
so all $|X_W| \geq \frac{k}{k-1}\dim W$ columns lie outside $\{0\}$, completing
the case $L = 0$. For any nonzero proper subspace $L \subsetneq W$, minimality
of $W$ gives $|X_L| < \frac{k}{k-1}\dim L$, and since $X_L \subseteq X_W$,
$$
        |\{x \in X_W : \mathbf{m}_x \notin L\}|
        = |X_W| - |X_L|
        > \frac{k}{k-1}(\dim W - \dim L).
$$
The vectors $(\mathbf{m}_x)_{x \in X_W}$ also span $W$, since if they spanned
only a proper subspace $W' \subsetneq W$ then $X_{W'} = X_W$ and $W'$ would
satisfy the defining inequality with smaller dimension, contradicting
minimality. Lemma~\ref{lem:complement-radohorn} therefore yields a partition
$X_W = P_1 \sqcup \cdots \sqcup P_k$ such that
\begin{equation}\label{eq:spans}
        \Span\{\mathbf{m}_x : x \in X_W \setminus P_j\} = W
        \quad \text{for every } j.
\end{equation}

We now show that~\eqref{eq:spans} gives the support condition for every
$\mathbf{c} \in \mathcal C$ that is not identically zero on $X_W$. Fix such
a $\mathbf{c}$, with coefficient vector $\boldsymbol{\beta}_{\mathbf{c}}$.
Since $\mathbf{c}$ is not identically zero on $X_W$, there exists $x_0 \in
X_W$ with $\boldsymbol{\beta}_{\mathbf{c}} \cdot \mathbf{m}_{x_0} \neq 0$,
so $\boldsymbol{\beta}_{\mathbf{c}}$ does not vanish on all of $W$. The set
$H = \{v \in W : \boldsymbol{\beta}_{\mathbf{c}} \cdot v = 0\}$ is therefore a
proper subspace of $W$. For each $j$,~\eqref{eq:spans} gives us that
$\{\mathbf{m}_x : x \in X_W \setminus P_j\}$ spans $W$, so these columns are
not all contained in $H$, and hence some $x_j \in X_W \setminus P_j$ satisfies
$\boldsymbol{\beta}_{\mathbf{c}} \cdot \mathbf{m}_{x_j} \neq 0$, meaning
$x_j \in \operatorname{supp} \mathbf{c} \setminus P_j$.

It remains to consider the vectors in $\mathcal C$ that vanish identically on
$X_W$. Their supports are contained in $Y = X^+ \setminus X_W$, and their
nonzero restrictions to $Y$ form a collection $\mathcal C_Y$. Fix $J \subseteq Y$.
An element of $\mathcal C_Y$ supported in $J$ is the restriction to $Y$ of
some $\mathbf{c}\in\mathcal C$ that vanishes on $X_W$. Since every vector in
$\mathcal C$ also vanishes on $X\setminus X^+$, and since the restriction of
$\mathbf{c}$ to $Y$ is supported in $J$, the original vector $\mathbf{c}$ is
supported in $J$ as a vector on $X$. Therefore
$\lambda_{\mathcal C_Y}(J) \leq \lambda_{\mathcal C}(J) \leq
\frac{k-1}{k}|J|$. Since $W \neq 0$ the set $X_W$ is nonempty, so $|Y| < |X|$,
and the induction hypothesis gives a partition $Y = Q_1 \sqcup \cdots \sqcup Q_k$
such that no vector in $\mathcal C_Y$ has support contained in a single $Q_j$.

Setting $Z_j = P_j \cup Q_j$ gives a partition of $X^+$. For any $\mathbf{c}
\in \mathcal C$ not identically zero on $X_W$, the preceding argument gives a
support point in $X_W \setminus P_j \subseteq Z_j^c$ for each $j$, so
$\operatorname{supp} \mathbf{c} \not\subseteq Z_j$. For any $\mathbf{c} \in
\mathcal C$ vanishing on $X_W$, its restriction lies in $\mathcal C_Y$, and
since $Z_j \cap Y = Q_j$ the inductive step gives $\operatorname{supp}
\mathbf{c} \not\subseteq Z_j$ for every $j$. Assigning the points of $X
\setminus X^+$ to any class completes the induction.
\end{proof}

\end{document}
